\epanalipsi{Β3}

\tmima{Λεξικό του κεφαλαίου}

\renewcommand{\arraystretch}{1.5}
\begin{tabular*}{\linewidth}{@{\extracolsep{\fill}} >{\raggedright\arraybackslash}p{0.28\linewidth} >{\raggedright\arraybackslash}p{0.62\linewidth}}
\toprule
\textbf{Λέξη ή σύμβολο} & \textbf{Με μία πρόταση}\\
\midrule
ισόπλευρο, ισοσκελές, σκαληνό & Τρεις, δύο ή καμία ίση πλευρά.\\
οξυγώνιο, ορθογώνιο, αμβλυγώνιο & Όλες οι γωνίες οξείες, μία ορθή, μία αμβλεία.\\
διάμεσος, ύψος, διχοτόμος & Προς το μέσο της απέναντι πλευράς, κάθετο στην απέναντι πλευρά, χωρίζει τη γωνία στη μέση.\\
$\hat{\mathrm{Α}}+\hat{\mathrm{Β}}+\hat{\Gamma}=180^\circ$ & Το άθροισμα των γωνιών τριγώνου.\\
εξωτερική γωνία & Ίση με το άθροισμα των δύο απέναντι εσωτερικών.\\
παραλληλόγραμμο & Απέναντι πλευρές παράλληλες, ίσες· διαγώνιες διχοτομούνται.\\
ορθογώνιο, ρόμβος, τετράγωνο & Ορθές γωνίες· ίσες πλευρές· και τα δύο μαζί.\\
τραπέζιο, ισοσκελές τραπέζιο & Μόνο δύο παράλληλες πλευρές· και ίσες οι άλλες δύο.\\
\bottomrule
\end{tabular*}

\tmima{Ασκήσεις επανάληψης}

\ex Τρίγωνο έχει πλευρές $9$, $9$ και $12$ cm. Ποιο είναι το είδος του και η περίμετρός του;

\ex Βρες την τρίτη γωνία τριγώνου με γωνίες $58^\circ$ και $64^\circ$. Τι είδους είναι;

\ex Ισοσκελές τρίγωνο έχει γωνία κορυφής $36^\circ$. Βρες τις γωνίες βάσης.

\ex Σε τρίγωνο η εξωτερική γωνία στο $\mathrm{Β}$ είναι $130^\circ$ και η $\hat{\mathrm{Α}}=45^\circ$. Βρες τις $\hat{\mathrm{Β}}$ και $\hat{\Gamma}$.

\ex Οι οξείες ενός ορθογωνίου τριγώνου είναι $x$ και $x+20^\circ$. Βρες τες.

\ex Πόσα ύψη έχει ένα τρίγωνο; Πού βρίσκονται σε αμβλυγώνιο;

\ex Βρες την περίμετρο ρόμβου με πλευρά $7{,}5$ cm και παραλληλογράμμου με πλευρές $9$ και $6$ cm.

\ex Ορθογώνιο $9\times12$ cm. Βρες τη διαγώνιο και την απόσταση του κέντρου συμμετρίας από μια κορυφή.

\ex Ρόμβος με διαγώνιες $14$ και $48$ cm. Βρες την πλευρά και την περίμετρό του.

\ex Παραλληλόγραμμο έχει γωνία $75^\circ$. Βρες τις άλλες.

\ex Ισοσκελές τραπέζιο έχει γωνίες βάσης $58^\circ$. Βρες τις άλλες.

\ex Πόσους άξονες συμμετρίας έχουν: ισόπλευρο τρίγωνο, ορθογώνιο, ρόμβος, τετράγωνο, ισοσκελές τραπέζιο;

\ex Η περίμετρος ισοσκελούς τραπεζίου είναι $46$ cm, οι βάσεις $16$ και $10$ cm. Βρες τις ίσες πλευρές.

\ex Ποιο παραλληλόγραμμο έχει ίσες και κάθετες διαγώνιες;

\ex \textbf{Ερωτήσεις αυτοαξιολόγησης (τρίγωνα).} Σωστό ή λάθος; α) Το ισόπλευρο έχει όλες τις γωνίες $60^\circ$. \quad β) Τρίγωνο έχει δύο ορθές γωνίες. \quad γ) Οι γωνίες βάσης ισοσκελούς είναι ίσες. \quad δ) Η εξωτερική γωνία ισούται με το άθροισμα των δύο απέναντι.

\ex \textbf{Ερωτήσεις αυτοαξιολόγησης (τετράπλευρα).} Σωστό ή λάθος; α) Κάθε ρόμβος είναι παραλληλόγραμμο. \quad β) Οι διαγώνιες ρόμβου είναι ίσες. \quad γ) Οι διαγώνιες ορθογωνίου είναι ίσες. \quad δ) Το τραπέζιο έχει τέσσερις παράλληλες πλευρές.

\tmima{Ο χάρτης της επανάληψης}

Μετά τις λύσεις, δες σε ποιες ασκήσεις δυσκολεύτηκες. Ο πίνακας σου λέει ποια παράγραφο να ξαναδιαβάσεις.

\renewcommand{\arraystretch}{1.5}
\begin{tabular*}{\linewidth}{@{\extracolsep{\fill}} >{\raggedright\arraybackslash}p{0.38\linewidth} >{\raggedright\arraybackslash}p{0.52\linewidth}}
\toprule
\textbf{Αν δυσκολεύτηκες στις} & \textbf{ξαναδιάβασε τις παραγράφους}\\
\midrule
ΕΒ3.1, ΕΒ3.6 & 3.1 (στοιχεία και είδη)\\
ΕΒ3.2, ΕΒ3.3, ΕΒ3.4, ΕΒ3.5 & 3.2 (άθροισμα γωνιών)\\
ΕΒ3.7 & 3.3 (είδη τετραπλεύρων)\\
ΕΒ3.8, ΕΒ3.9, ΕΒ3.10, ΕΒ3.11, ΕΒ3.13 & 3.4 (ιδιότητες)\\
ΕΒ3.12, ΕΒ3.14 & 3.4 και Β2 (συμμετρία)\\
ΕΒ3.15, ΕΒ3.16 & 3.1 έως 3.4\\
\bottomrule
\end{tabular*}
